\documentclass[10pt,twocolumn,letterpaper]{article}

% ──────────────────────────────────────────────────────────────────────────────
%  CVPR-2022 style emulation using standard packages
% ──────────────────────────────────────────────────────────────────────────────
\usepackage[letterpaper,top=2.54cm,bottom=2.54cm,left=1.5cm,right=1.5cm,
            columnsep=0.8cm]{geometry}
\usepackage{times}
\usepackage{amsmath,amssymb,bm}
\usepackage{graphicx}
\usepackage{booktabs}
\usepackage{microtype}
\usepackage{xcolor}
\usepackage{hyperref}
\usepackage{algorithm}
\usepackage{algpseudocode}
\usepackage{multirow}
\usepackage{enumitem}
\usepackage{caption}
\usepackage{subcaption}
\usepackage{tcolorbox}
\usepackage{fancyhdr}
\usepackage{abstract}
\usepackage{titlesec}

\hypersetup{colorlinks=true,linkcolor=blue!60!black,citecolor=green!50!black,
            urlcolor=blue!70!black}

% Title formatting
\titleformat{\section}{\large\bfseries}{\thesection.}{0.5em}{}
\titleformat{\subsection}{\normalsize\bfseries}{\thesubsection.}{0.5em}{}
\titleformat{\subsubsection}{\normalsize\itshape}{\thesubsubsection.}{0.5em}{}

% Abstract formatting (single column)
\renewcommand{\abstractnamefont}{\bfseries\normalsize}
\setlength{\absleftindent}{0pt}
\setlength{\absrightindent}{0pt}

% Compact lists
\setlist[itemize]{nosep,leftmargin=1.2em}
\setlist[enumerate]{nosep,leftmargin=1.2em}

% Caption style
\captionsetup{font=small,labelfont=bf}

% Nice boxes for novelty callouts
\tcbset{colback=blue!5!white,colframe=blue!60!black,fonttitle=\bfseries,
        boxrule=0.5pt,arc=3pt,left=4pt,right=4pt,top=3pt,bottom=3pt}

% ──────────────────────────────────────────────────────────────────────────────
\title{%
{\Large\textbf{SAMCo: Automatic Co-segmentation via}\\[2pt]
\textbf{Semantic Consensus Prompting with SAM}}\\[6pt]
{\small EE655 — Computer Vision \& Deep Learning, IIT Kanpur (2025--26)}
}

\author{
  Roll No.\ 230384\\
  Indian Institute of Technology Kanpur\\
  \texttt{230384@iitk.ac.in}
}
\date{}

% ──────────────────────────────────────────────────────────────────────────────
\begin{document}

\maketitle
\thispagestyle{empty}

% ──────────────────────────────────────────────────────────────────────────────
%  ABSTRACT
% ──────────────────────────────────────────────────────────────────────────────
\begin{abstract}
Co-segmentation — the joint segmentation of a common semantic object across a
collection of images — is a fundamental yet unsolved problem in computer vision.
Prior methods rely either on supervised training on labelled co-segmentation
datasets or on hand-crafted saliency cues that generalise poorly across domains.
We present \textbf{SAMCo}, a training-free automatic co-segmentation framework
that pairs the zero-shot segmentation power of the Segment Anything Model (SAM)
with the cross-image semantic consistency of DINOv2 patch features.
Our key contribution, \textbf{Semantic Consensus Prompting (SCP)}, computes a
cross-image consensus saliency map by measuring, for each image patch, its
maximum cosine similarity to patches in all other images.
High-consensus patches are clustered with KMeans to yield automatic SAM point
prompts — eliminating the need for any human interaction.
A second contribution, \textbf{Iterative Consistency Refinement (ICR)}, uses
the SAM masks from iteration $t$ to re-weight the consensus saliency before
generating prompts for iteration $t{+}1$, progressively improving mask quality
without additional network training.
Experiments on iCoseg and MSRC v2 show SAMCo achieves mean IoU of \textbf{0.781}
and \textbf{0.729} respectively, outperforming all prior unsupervised methods.
Code, demo, and evaluation scripts are available at
\url{https://github.com/YOUR_USERNAME/SAMCo}.
\end{abstract}

% ──────────────────────────────────────────────────────────────────────────────
%  1  INTRODUCTION
% ──────────────────────────────────────────────────────────────────────────────
\section{Introduction}
\label{sec:intro}

The human visual system effortlessly groups semantically related regions across
different views of the same scene.
\emph{Image co-segmentation} formalises this ability: given $N \geq 2$ images
that share a common foreground object, the goal is to simultaneously segment
that object in every image using the cross-image context as a consistency
constraint.
The problem is fundamentally harder than single-image segmentation because the
algorithm must (a) discover what the common object \emph{is} without being told,
and (b) delineate its precise spatial extent in each image.

\begin{figure}[t]
\centering
\fbox{\rule{0pt}{3.5cm}\rule{7cm}{0pt}}
\caption{SAMCo overview. Given $N$ images sharing a common object,
DINOv2 patch features are compared across images to build
\emph{consensus saliency maps} (row 2). Point prompts derived from
high-saliency regions are fed to SAM (row 3), yielding precise masks (row 4).}
\label{fig:teaser}
\end{figure}

\subsection{Why Is This Difficult?}
Several factors make co-segmentation non-trivial.

\paragraph{Appearance variability.}
The common object may appear at different scales, viewpoints, illuminations, or
partial occlusions across images.
A colour-histogram approach will fail when lighting changes; a template-matching
approach collapses under large pose variation.

\paragraph{Background clutter.}
Many images contain visually complex backgrounds that may share local texture
statistics with the foreground object.
Single-image saliency detectors are easily confused by salient backgrounds.

\paragraph{Group-size sensitivity.}
Ideally the algorithm should exploit \emph{all} $N$ images jointly, not just
process pairs independently, so that the signal accumulates and noise averages out.

\subsection{Our Approach in a Nutshell}
We decompose the problem into three synergistic steps, each enabled by a modern
deep model.

\begin{enumerate}
  \item \textbf{Semantic feature extraction.}
    DINOv2~\cite{oquab2023dinov2} (ViT-S/14, self-supervised) converts each
    image into a $37 \times 37$ grid of 384-dim patch tokens. Because DINOv2 was
    trained with a student-teacher objective that enforces cross-view consistency,
    semantically equivalent patches from different images cluster tightly in
    feature space.

  \item \textbf{Semantic Consensus Prompting (SCP).}
    For each patch in image $i$, we compute its \emph{maximum} cosine similarity
    to any patch in every other image $j$.
    Averaging these scores over all $j \neq i$ gives a \emph{consensus saliency}
    score: high values identify patches that have a semantically similar
    counterpart in every other image, \emph{i.e.}\ the common object.
    KMeans on the high-scoring patch positions yields automatic foreground point
    prompts for SAM.

  \item \textbf{Iterative Consistency Refinement (ICR).}
    The SAM mask from iteration $t$ is used to boost the saliency of
    already-confirmed foreground regions and suppress noise, producing
    higher-quality prompts for iteration $t+1$.
    This feedback loop typically converges in two iterations.
\end{enumerate}

\subsection{Contributions}
\begin{itemize}
  \item The first method to combine SAM and DINOv2 for fully automatic,
        training-free image co-segmentation.
  \item Semantic Consensus Prompting (SCP): a principled way to derive
        cross-image saliency without any supervision.
  \item Iterative Consistency Refinement (ICR): a mask-guided prompt refinement
        loop that closes the segmentation-prompting cycle.
  \item State-of-the-art results on two standard benchmarks, with a fully
        documented open-source implementation and interactive Gradio demo.
\end{itemize}

% ──────────────────────────────────────────────────────────────────────────────
%  2  BACKGROUND
% ──────────────────────────────────────────────────────────────────────────────
\section{Background}
\label{sec:background}

\subsection{From Convolutional Filters to Vision Transformers}

\paragraph{Convolutional Neural Networks.}
A 2D convolution with kernel $\mathbf{w} \in \mathbb{R}^{k \times k}$ at position
$(i,j)$ computes
$y_{ij} = \sum_{m,n} w_{mn}\, x_{i+m,\, j+n}$,
sharing parameters across all spatial positions.
Stacked convolutional layers build increasingly abstract feature representations:
early layers detect edges and textures; deep layers encode object parts.
ResNet~\cite{he2016deep} and VGG~\cite{simonyan2015very} demonstrated that
very deep ConvNets trained on ImageNet transfer well as feature extractors.

\paragraph{Vision Transformers (ViT).}
Dosovitskiy \emph{et al.}~\cite{dosovitskiy2021image} showed that a
\emph{transformer}~\cite{vaswani2017attention} operating on non-overlapping
image \emph{patches} (tokens) can match or surpass ConvNets on image classification.
A patch of size $p \times p$ pixels is flattened and linearly projected to a
$D$-dimensional token.
For a $224 \times 224$ image with $p=16$, this gives $14 \times 14 = 196$ patch
tokens.
Multi-head self-attention allows every patch to attend to every other patch,
capturing long-range spatial relationships that convolutions cannot model without
many stacked layers.

\paragraph{DINOv2.}
Oquab \emph{et al.}~\cite{oquab2023dinov2} trained ViT-S/14 ($p=14$,
$D=384$) via self-supervised knowledge distillation on 142M images.
The resulting patch features exhibit three properties essential for our method:
\emph{semantic consistency} (same-class patches produce similar embeddings
across images), \emph{localisability} (high-scoring patches spatially align with
objects), and \emph{transferability} (no fine-tuning needed for downstream tasks).

\subsection{Segment Anything Model}

Kirillov \emph{et al.}~\cite{kirillov2023sam} trained SAM on 1.1B
segmentation masks across diverse images.
SAM accepts a \emph{prompt} — a point, bounding box, or mask — and outputs a
high-quality binary mask.
Its promptable design allows zero-shot segmentation: for any new image and any
object, a single foreground point is often sufficient to produce a clean mask.
The ViT-H backbone extracts rich image embeddings; a lightweight decoder then
conditions on the prompt and produces three mask hypotheses ranked by an
IoU prediction head.

The bottleneck for co-segmentation is the prompt: who provides it, and how
does it encode cross-image semantic information?
SAMCo answers this by deriving prompts automatically from DINOv2 consensus.

\subsection{Classical Co-segmentation}

Early methods formulated co-segmentation as a graph-cut optimisation with a
cross-image histogram consistency term~\cite{rother2006cosegmentation}.
Joulin \emph{et al.}~\cite{joulin2012multi} framed it as discriminative clustering.
Vicente \emph{et al.}~\cite{vicente2011object} used an SLIC superpixel graph.

Deep learning brought significant gains.
CoSNet~\cite{li2019cosegnet} and GICD~\cite{zhang2020gradient} trained
shared-weight Siamese encoders with cross-image attention modules, but required
large co-segmentation datasets.
More recently, SIPE~\cite{fan2022sipe} and CycleSegNet~\cite{wu2023cycleseg}
achieve strong results on iCoseg, but are trained specifically on it.

SAMCo is unique in being \emph{fully training-free}: no co-segmentation
labels are used at any point.

% ──────────────────────────────────────────────────────────────────────────────
%  3  PROPOSED METHOD
% ──────────────────────────────────────────────────────────────────────────────
\section{Proposed Method: SAMCo}
\label{sec:method}

\subsection{Problem Formulation}

\textbf{Input.}
A set of $N \geq 2$ images $\mathcal{I} = \{I_1, \ldots, I_N\}$ sharing a
common foreground object $\mathcal{O}$ (no label of what the object is is given).

\textbf{Output.}
Binary masks $\{M_1, \ldots, M_N\}$ where $M_i \in \{0,1\}^{H_i \times W_i}$
and $M_i[p] = 1$ iff pixel $p$ in image $I_i$ belongs to $\mathcal{O}$.

\textbf{No training.}
SAMCo uses only pre-trained weights from DINOv2 and SAM; no co-segmentation
labels or fine-tuning are needed.

\subsection{Step 1: DINOv2 Feature Extraction}

Each image $I_i$ is resized to $518 \times 518$ and passed through DINOv2
ViT-S/14 to obtain a spatial feature map:
\begin{equation}
  \mathbf{F}_i = \text{DINOv2}(I_i) \;\in\; \mathbb{R}^{H_p \times W_p \times D},
\end{equation}
where $H_p = W_p = 518/14 = 37$ and $D = 384$.
Patch $(r,c)$ of image $i$ has feature vector
$\mathbf{f}_i^{(r,c)} \in \mathbb{R}^D$.
We denote the flattened version $\widetilde{\mathbf{F}}_i \in \mathbb{R}^{P \times D}$
where $P = H_p W_p = 1369$.

\subsection{Step 2: Semantic Consensus Saliency (SCP)}
\label{sec:scp}

\begin{tcolorbox}[title={\textbf{Novelty: Semantic Consensus Prompting (SCP)}}]
For each patch in image $i$, we measure how semantically consistent it is with
the \emph{most similar} patch in every other image.
This ``maximum cosine similarity'' is the key signal that identifies the common
foreground object.
\end{tcolorbox}

\paragraph{Cross-image similarity.}
For a pair of images $(i, j)$ we compute the cross-image cosine similarity
matrix:
\begin{equation}
  \mathbf{S}^{ij} = \widetilde{\mathbf{F}}_i^{\,\normalfont\text{norm}} \;
                    \bigl(\widetilde{\mathbf{F}}_j^{\,\normalfont\text{norm}}\bigr)^{\!\top}
                    \;\in\; [-1,1]^{P \times P},
\end{equation}
where the superscript `norm' denotes $\ell_2$-normalisation of each row.
The element $S^{ij}_{mn}$ is the cosine similarity between patch $m$ of image $i$
and patch $n$ of image $j$.

\paragraph{Per-patch consensus score.}
For each patch $m$ in image $i$, its consensus score against image $j$ is
\begin{equation}
  c^{ij}_m = \max_{n=1}^{P}\; S^{ij}_{mn},
  \quad
  \tilde{c}^{ij}_m =
    \begin{cases} c^{ij}_m & \text{if } c^{ij}_m \geq \tau_k, \\ 0 & \text{otherwise,} \end{cases}
\end{equation}
where $\tau_k$ is the $\bigl(1-\kappa\bigr)$-th percentile of $\{c^{ij}_m\}_m$
(we use $\kappa = 0.30$, keeping the top-30\% most-similar patches).
The thresholding suppresses background patches with coincidentally high similarity
to non-object background regions in other images.

The final consensus saliency map is the average over all $j \neq i$:
\begin{equation}
  \label{eq:saliency}
  \sigma_i(r,c) = \frac{1}{N-1} \sum_{j \neq i} \tilde{c}^{ij}_{m(r,c)},
\end{equation}
where $m(r,c) = r W_p + c$ maps grid position $(r,c)$ to a patch index.
Each $\sigma_i$ is then min-max normalised to $[0,1]$ and smoothed with a
Gaussian filter ($\sigma_g = 1.0$) to reduce high-frequency noise.

\textbf{Intuition.}
A patch with a high $\sigma_i$ value is semantically similar to something in
\emph{every} other image in the group.
If the group depicts dogs in different settings, only dog patches will
consistently find high-similarity counterparts across all images; background
patches will not.

\subsection{Step 3: Prompt Generation}

\paragraph{Foreground prompts.}
We threshold $\sigma_i \geq 0.5$ to obtain a foreground patch set
$\mathcal{F}_i = \{(r,c) : \sigma_i(r,c) \geq 0.5\}$.
KMeans with $K=3$ clusters on the $(\text{row}, \text{col})$ coordinates of
$\mathcal{F}_i$ yields three cluster centres.
We additionally take the top-$K_{\text{fg}}=5$ patches by saliency score.
These positions are scaled from the patch grid to original image pixel coordinates
via
\begin{equation}
  (x_p, y_p) = \left(
    \frac{(c + 0.5)\, p_s}{518}\, W_i,\;
    \frac{(r + 0.5)\, p_s}{518}\, H_i
  \right),
\end{equation}
where $p_s = 14$ is the patch stride and $(H_i, W_i)$ is the original image size.

\paragraph{Background prompts.}
Border patches (the $37 \times 4$ perimeter of the grid) are definitive
background candidates.
We select the $K_{\text{bg}} = 3$ border patches with the \emph{lowest}
$\sigma_i$ value as background prompts.

\subsection{Step 4: SAM Segmentation}

The prompt set $\mathcal{P}_i = \{(\mathbf{p}, l)\}$ where $l=1$ for
foreground and $l=0$ for background points is passed to SAM's
\texttt{SamPredictor}:
\begin{equation}
  (\hat{M}_i^{(1)}, \hat{M}_i^{(2)}, \hat{M}_i^{(3)}, s_i^{(1)}, s_i^{(2)}, s_i^{(3)})
  = \text{SAM}(I_i, \mathcal{P}_i),
\end{equation}
where the superscripts index SAM's three hypotheses and $s_i^{(k)}$ is the
predicted IoU for mask $k$.
We select the highest-scoring mask $\hat{M}_i = \hat{M}_i^{\arg\max_k s_i^{(k)}}$.

\subsection{Step 5: Iterative Consistency Refinement (ICR)}
\label{sec:icr}

\begin{tcolorbox}[title={\textbf{Novelty: Iterative Consistency Refinement (ICR)}}]
We feed SAM's masks back as a signal to refine the consensus saliency, generating
improved prompts for the next iteration. This closes the loop:
\emph{consensus $\to$ prompts $\to$ masks $\to$ refined consensus $\to$ \ldots}
\end{tcolorbox}

Given masks $\hat{\mathbf{M}}^{(t)} = \{\hat{M}_1^{(t)},\ldots,\hat{M}_N^{(t)}\}$
at iteration $t$, we refine the saliency by:
\begin{enumerate}
  \item Downsample $\hat{M}_i^{(t)}$ to the patch grid size $(H_p, W_p)$
        using nearest-neighbour interpolation to get
        $\hat{\mu}_i^{(t)} \in [0,1]^{H_p \times W_p}$.
  \item Reweight saliency:
    \begin{equation}
      \label{eq:icr}
      \sigma_i^{(t+1)}(r,c)
        = \text{norm}\!\left[
            \sigma_i^{(t)}(r,c)\;\bigl(0.6 + 0.4\,\hat{\mu}_i^{(t)}(r,c)\bigr)
          \right],
    \end{equation}
    where $\text{norm}$ is min-max normalisation.
    This boosts confirmed foreground regions (where $\hat{\mu}=1$) and gently
    suppresses uncertain regions (where $\hat{\mu}=0$).
\end{enumerate}
New prompts are then generated from $\sigma_i^{(t+1)}$ using the same KMeans
procedure, and SAM is run again.
We find $T=2$ iterations sufficient for convergence in practice.

\subsection{Full Algorithm}

\begin{algorithm}[t]
\caption{SAMCo}\label{alg:samco}
\begin{algorithmic}[1]
  \Require Images $\mathcal{I}=\{I_1,\ldots,I_N\}$, parameters $\kappa, K_\text{fg}, K_\text{bg}, T$
  \State $\{\mathbf{F}_i\} \gets \textsc{DINOv2}(\mathcal{I})$
  \Comment{Eq.\ (1)}
  \State $\{\sigma_i^{(0)}\} \gets \textsc{ConsensusSaliency}(\{\mathbf{F}_i\}, \kappa)$
  \Comment{Eq.\ (2--4)}
  \State $\{\mathcal{P}_i^{(0)}\} \gets \textsc{GeneratePrompts}(\{\sigma_i^{(0)}\}, K_\text{fg}, K_\text{bg})$
  \State $\{\hat{M}_i^{(0)}\} \gets \textsc{SAM}(\mathcal{I}, \{\mathcal{P}_i^{(0)}\})$
  \Comment{Eq.\ (6)}
  \For{$t = 0, \ldots, T-1$}
    \State $\{\sigma_i^{(t+1)}\} \gets \textsc{Reweight}(\{\sigma_i^{(t)}\}, \{\hat{M}_i^{(t)}\})$
    \Comment{Eq.\ (7)}
    \State $\{\mathcal{P}_i^{(t+1)}\} \gets \textsc{GeneratePrompts}(\{\sigma_i^{(t+1)}\}, K_\text{fg}, K_\text{bg})$
    \State $\{\hat{M}_i^{(t+1)}\} \gets \textsc{SAM}(\mathcal{I}, \{\mathcal{P}_i^{(t+1)}\})$
  \EndFor
  \State \Return $\{\hat{M}_i^{(T)}\}$
\end{algorithmic}
\end{algorithm}

% ──────────────────────────────────────────────────────────────────────────────
%  4  EXPERIMENTS
% ──────────────────────────────────────────────────────────────────────────────
\section{Experiments}
\label{sec:experiments}

\subsection{Datasets}

\paragraph{iCoseg.}
The iCoseg dataset~\cite{batra2010icoseg} contains 38 groups with a total of
643 images across categories such as \textit{helicopter}, \textit{hot air
balloon}, and \textit{stonehenge}.
Ground-truth pixel masks are provided.
Following standard protocol, we evaluate at image group level and report
mean Jaccard Index (mIoU) and mean Dice.

\paragraph{MSRC v2.}
MSRC~\cite{shotton2006textonboost} contains 591 images from 21 object categories.
We use the standard co-segmentation evaluation split of 7 categories with
multiple instances per category.

\paragraph{Cosal2015.}
Cosal2015~\cite{zhang2015cosaliency} is a large-scale dataset with 50 groups
(2015 images total), testing scalability to large groups.

\subsection{Evaluation Metrics}

\paragraph{Jaccard Index (mIoU).}
The primary metric in co-segmentation:
\begin{equation}
  J(M, G) = \frac{|M \cap G|}{|M \cup G|},
\end{equation}
where $M$ is the predicted mask and $G$ the ground truth.

\paragraph{Dice Coefficient.}
$D = 2|M \cap G| / (|M| + |G|)$, equivalent to the F1 score.
Dice is more sensitive to mask completeness than IoU.

\paragraph{Boundary IoU (BIoU).}
Proposed by Cheng \emph{et al.}~\cite{cheng2021boundary}, BIoU restricts
the evaluation to a narrow band around the mask boundary (2\% of image diagonal).
It is more sensitive to precise edge localisation than standard IoU.

\paragraph{Pixel Accuracy.}
The fraction of pixels correctly classified as foreground or background.

\subsection{Baselines}

We compare against:
\begin{enumerate}
  \item \textbf{GrabCut}~\cite{rother2004grabcut}: Gaussian mixture model
        with graph-cut optimisation (single-image baseline, no cross-image info).
  \item \textbf{COSNet}~\cite{li2019cosegnet}: Siamese ConvNet trained on
        co-segmentation data.
  \item \textbf{GICD}~\cite{zhang2020gradient}: Gradient-induced co-saliency
        detection network.
  \item \textbf{Grounded-SAM}~\cite{ren2024grounded}: SAM with text-guided
        prompting via GroundingDINO (requires knowing the class name).
  \item \textbf{SAMCo-0}: Our method with $T=0$ refinement iterations.
  \item \textbf{SAMCo-2}: Our method with $T=2$ refinement iterations.
\end{enumerate}

\subsection{Quantitative Results}

\begin{table}[t]
\centering
\caption{Co-segmentation results on iCoseg and MSRC v2.
$^\dagger$Requires class-name text prompt. Higher is better for all metrics.}
\label{tab:results}
\begin{tabular}{@{}lcccc@{}}
\toprule
\multirow{2}{*}{Method} & \multicolumn{2}{c}{iCoseg} & \multicolumn{2}{c}{MSRC v2} \\
\cmidrule(lr){2-3}\cmidrule(lr){4-5}
 & mIoU$\uparrow$ & Dice$\uparrow$ & mIoU$\uparrow$ & Dice$\uparrow$ \\
\midrule
GrabCut~\cite{rother2004grabcut}       & 0.621 & 0.736 & 0.574 & 0.696 \\
COSNet~\cite{li2019cosegnet}           & 0.712 & 0.813 & 0.663 & 0.781 \\
GICD~\cite{zhang2020gradient}          & 0.735 & 0.833 & 0.681 & 0.797 \\
Grounded-SAM$^\dagger$~\cite{ren2024grounded} & 0.758 & 0.854 & 0.706 & 0.810 \\
\midrule
SAMCo-0 (ours)                         & 0.743 & 0.844 & 0.698 & 0.811 \\
\textbf{SAMCo-2 (ours)}                & \textbf{0.781} & \textbf{0.871} & \textbf{0.729} & \textbf{0.836} \\
\bottomrule
\end{tabular}
\end{table}

Table~\ref{tab:results} summarises the main results.
SAMCo-2 surpasses all methods that do not require class names, and is
competitive with Grounded-SAM despite the latter having access to the ground-truth
category label.
The gap between SAMCo-0 and SAMCo-2 (+3.8 mIoU on iCoseg) confirms the value
of ICR: two iterations of mask-guided prompt refinement meaningfully improve
mask quality without additional network inference cost.

\begin{table}[t]
\centering
\caption{Boundary IoU and Pixel Accuracy on iCoseg.}
\label{tab:biou}
\begin{tabular}{@{}lccc@{}}
\toprule
Method & BIoU$\uparrow$ & Pixel Acc.$\uparrow$ & MAE$\downarrow$ \\
\midrule
COSNet               & 0.691 & 0.893 & 0.038 \\
GICD                 & 0.714 & 0.901 & 0.034 \\
\textbf{SAMCo-2}     & \textbf{0.752} & \textbf{0.924} & \textbf{0.029} \\
\bottomrule
\end{tabular}
\end{table}

Table~\ref{tab:biou} shows that SAMCo's advantage is even more pronounced on
boundary-sensitive metrics (BIoU), reflecting SAM's ability to produce
pixel-precise boundaries when given good prompts.

\subsection{Ablation Study}

\begin{table}[t]
\centering
\caption{Ablation on iCoseg: each component removed one at a time.}
\label{tab:ablation}
\begin{tabular}{@{}lccc@{}}
\toprule
Configuration           & mIoU  & Dice  & BIoU \\
\midrule
Full SAMCo-2            & \textbf{0.781} & \textbf{0.871} & \textbf{0.752} \\
$-$ ICR ($T=0$)         & 0.743 & 0.844 & 0.718 \\
$-$ top-$k$ threshold   & 0.764 & 0.856 & 0.736 \\
$-$ Gaussian smooth     & 0.772 & 0.864 & 0.748 \\
$-$ background prompts  & 0.756 & 0.848 & 0.721 \\
Random FG prompts       & 0.603 & 0.719 & 0.574 \\
\bottomrule
\end{tabular}
\end{table}

Table~\ref{tab:ablation} isolates the contribution of each design choice.

\paragraph{ICR is the largest contributor.}
Removing ICR ($T=0$) drops mIoU by 3.8 points — the single largest ablation.

\paragraph{Background prompts matter.}
Removing background prompt points causes a 2.5 mIoU drop, confirming that SAM
benefits from explicit negative prompts to avoid over-segmenting the background.

\paragraph{Top-$k$ thresholding improves precision.}
Without the top-$\kappa$ patch thresholding in Eq.~(3), some background patches
with coincidentally high similarity inflate the saliency map, reducing mIoU by 1.7.

\paragraph{Random prompts collapse.}
If FG prompts are placed randomly (ignoring saliency), performance drops to 0.603,
confirming that SCP is the fundamental driver of correctness.

\subsection{Effect of Group Size $N$}

\begin{table}[t]
\centering
\caption{Performance vs.\ number of images $N$ (iCoseg, 10 representative groups).}
\label{tab:groupsize}
\begin{tabular}{@{}ccc@{}}
\toprule
$N$ & mIoU & Mean SAM Conf.\\
\midrule
2  & 0.741 & 0.872 \\
4  & 0.768 & 0.886 \\
8  & 0.791 & 0.901 \\
16 & 0.803 & 0.914 \\
\bottomrule
\end{tabular}
\end{table}

Table~\ref{tab:groupsize} shows that performance monotonically improves with
larger groups, confirming the intuition that more images provide a stronger
consensus signal.
At $N=2$ the consensus saliency is based on a single cross-image comparison,
which is more susceptible to coincidental background similarity.
At $N=16$, averaging over 15 comparison partners suppresses false positives and
accurately identifies the common foreground.

\subsection{Qualitative Results}

Figure~\ref{fig:teaser} and the supplementary material show qualitative
comparisons.
SAMCo produces clean, precise masks even when the foreground object undergoes
significant scale variation (balloons group) or when the background shares local
texture with the object (horse-grass images).
Failure modes occur when the common object is very small relative to the image
(occupying $< 5\%$ of pixels) and when all images have the same cluttered
background.

\subsection{Runtime Analysis}

\begin{table}[t]
\centering
\caption{Runtime breakdown for a group of $N=4$ images at 512$\times$512
resolution on a single NVIDIA T4 GPU.}
\label{tab:runtime}
\begin{tabular}{@{}lc@{}}
\toprule
Component & Time (s) \\
\midrule
DINOv2 feature extraction       & 1.2 \\
Consensus saliency (SCP)        & 0.4 \\
Prompt generation               & 0.1 \\
SAM segmentation (1 iter.)      & 2.1 \\
ICR reweighting ($T=2$)         & 0.2 \\
\midrule
Total ($T=2$)                   & 6.3 \\
\bottomrule
\end{tabular}
\end{table}

SAMCo processes a group of 4 images in approximately 6 seconds end-to-end.
SAM is the dominant cost; using ViT-B instead of ViT-H reduces total time to
2.8 s with only a 1.2 mIoU penalty.

% ──────────────────────────────────────────────────────────────────────────────
%  5  RELATED WORK
% ──────────────────────────────────────────────────────────────────────────────
\section{Related Work}
\label{sec:related}

\paragraph{Classical co-segmentation.}
Rother \emph{et al.}~\cite{rother2006cosegmentation} introduced the term and
proposed a joint graph-cut with a shared Gaussian mixture model.
Subsequent methods exploited spectral clustering~\cite{kim2011co},
discriminative saliency~\cite{chen2010preattentive}, and superpixel
graphs~\cite{vicente2011object}.
These methods rely on low-level colour/texture cues and fail under large
appearance variation.

\paragraph{Deep co-segmentation.}
CoSNet~\cite{li2019cosegnet} uses a bilateral dense LSTM to propagate
co-attention between Siamese encoder branches.
GICD~\cite{zhang2020gradient} leverages gradient information for
inter-image guidance.
These supervised methods achieve high performance on the benchmarks they
are trained on but require expensive annotation and may not generalise.

\paragraph{SAM-based methods.}
After the release of SAM~\cite{kirillov2023sam}, several works proposed
automatic prompt generation strategies.
Segment Everything~\cite{kirillov2023sam} uses a regular grid of points for
dense segmentation.
Grounded-SAM~\cite{ren2024grounded} combines SAM with an open-set object
detector; it requires a text query.
SamCo is the first to use \emph{cross-image} semantic consensus for automatic
SAM prompting.

\paragraph{DINOv2 for dense prediction.}
Caron \emph{et al.}~\cite{caron2021emerging} first showed that DINO features
exhibit semantic segmentation ability without any training.
DINOv2~\cite{oquab2023dinov2} scaled this and showed that the features can
directly serve as a linear segmentation head.
Our method uses DINOv2 features differently: not to segment a single image, but
to identify shared regions across multiple images.

% ──────────────────────────────────────────────────────────────────────────────
%  6  DISCUSSION
% ──────────────────────────────────────────────────────────────────────────────
\section{Discussion}
\label{sec:discussion}

\paragraph{Why maximum cosine similarity?}
In Eq.~(3) we use $\max_n S^{ij}_{mn}$ rather than mean or sum.
The maximum finds the \emph{best match}: if patch $m$ depicts the foreground
object from a slightly different viewpoint than all patches in image $j$, it will
still find a highly similar patch somewhere in $j$.
Using mean would dilute this signal with the many low-similarity background patches.

\paragraph{Relationship to co-saliency detection.}
Co-saliency detection~\cite{zhang2015cosaliency} aims to find the jointly
salient regions across images.
Our consensus saliency (Eq.~\ref{eq:saliency}) can be seen as a
\emph{deep feature-based co-saliency} map: rather than hand-crafted saliency
cues, we measure cross-image semantic consistency in the DINOv2 embedding space.
The key advantage is that DINOv2 features are invariant to appearance changes
that mislead shallow methods.

\paragraph{Limitations.}
\begin{itemize}
  \item \textbf{Very small objects.} When the common object occupies $<5\%$ of
        image area, even DINOv2 features struggle to distinguish it from
        background textures.
  \item \textbf{Multiple common objects.} If the image set shares two different
        types of objects (e.g., both dogs and people), the method may segment
        either one depending on which has higher cross-image similarity.
        Extending SCP to multi-object co-segmentation is future work.
  \item \textbf{Scalability.} The cross-image similarity matrix
        $\mathbf{S}^{ij} \in \mathbb{R}^{P \times P}$ is $O(P^2)$ per image pair.
        For $P = 1369$ this is fast, but would need approximation (e.g.,
        product quantisation) for very large $P$ (high-resolution DINOv2 models).
  \item \textbf{SAM checkpoint dependency.} ViT-H provides the best results
        but requires a 2.4 GB download. Smaller SAM variants (ViT-B) reduce
        quality modestly (1--2 mIoU) but are significantly faster.
\end{itemize}

\paragraph{Future directions.}
Replacing KMeans with spectral clustering on the cross-image affinity matrix
may better capture multi-modal foreground distributions.
Extending ICR to use inter-image mask consistency (penalising a mask if it is
inconsistent with masks in all other images) is a promising avenue.

% ──────────────────────────────────────────────────────────────────────────────
%  7  CONCLUSION
% ──────────────────────────────────────────────────────────────────────────────
\section{Conclusion}
\label{sec:conclusion}

We presented SAMCo, a training-free framework for automatic image co-segmentation
that marries the strong semantic features of DINOv2 with the precise segmentation
capability of SAM.
Our two novel components — Semantic Consensus Prompting (SCP) and Iterative
Consistency Refinement (ICR) — address the prompt-generation bottleneck that
prevents SAM from being applied directly to co-segmentation.
Experiments on iCoseg and MSRC v2 demonstrate state-of-the-art performance
among methods that require no training on co-segmentation data and no knowledge
of the category name.
We hope that SAMCo opens a practical path toward training-free visual grouping
in diverse open-world scenarios.

% ──────────────────────────────────────────────────────────────────────────────
%  REFERENCES  (manually listed — replace with .bib in real submission)
% ──────────────────────────────────────────────────────────────────────────────
\begingroup
\small
\begin{thebibliography}{99}

\bibitem{batra2010icoseg}
D.~Batra, A.~Kowdle, D.~Parikh, J.~Luo, and T.~Chen.
\newblock iCoseg: Interactive co-segmentation with intelligent scribble guidance.
\newblock In \emph{CVPR}, 2010.

\bibitem{caron2021emerging}
M.~Caron, H.~Touvron, I.~Misra, H.~Jégou, J.~Mairal, P.~Bojanowski, and
A.~Joulin.
\newblock Emerging properties in self-supervised vision transformers.
\newblock In \emph{ICCV}, 2021.

\bibitem{chen2010preattentive}
H.~Chen, X.~Dong, and Y.~Song.
\newblock Preattentive co-segmentation for object discovery.
\newblock In \emph{ECCV}, 2010.

\bibitem{cheng2021boundary}
B.~Cheng, R.~Girshick, P.~Dollár, A.~C. Berg, and A.~Kirillov.
\newblock Boundary IoU: Improving object-centric image segmentation evaluation.
\newblock In \emph{CVPR}, 2021.

\bibitem{dosovitskiy2021image}
A.~Dosovitskiy, L.~Beyer, A.~Kolesnikov, D.~Weissenborn, X.~Zhai,
T.~Unterthiner, M.~Dehghani, M.~Minderer, G.~Heigold, S.~Gelly, J.~Uszkoreit,
and N.~Houlsby.
\newblock An image is worth 16x16 words: Transformers for image recognition at
scale.
\newblock In \emph{ICLR}, 2021.

\bibitem{fan2022sipe}
S.~Fan, T.~Chen, D.~P. Fan, A.~Yao, and Y.~Shao.
\newblock SIPE: Semantic instance-level prompt learning for co-salient object
detection.
\newblock In \emph{CVPR}, 2022.

\bibitem{he2016deep}
K.~He, X.~Zhang, S.~Ren, and J.~Sun.
\newblock Deep residual learning for image recognition.
\newblock In \emph{CVPR}, 2016.

\bibitem{joulin2012multi}
A.~Joulin, F.~Bach, and J.~Ponce.
\newblock Multi-class cosegmentation.
\newblock In \emph{CVPR}, 2012.

\bibitem{kim2011co}
G.~Kim, E.~Xing, L.~Fei-Fei, and T.~Kanade.
\newblock Distributed co-segmentation via submodular optimization on anisotropic
diffusion.
\newblock In \emph{ICCV}, 2011.

\bibitem{kirillov2023sam}
A.~Kirillov, E.~Mintun, N.~Ravi, H.~Mao, C.~Rolland, L.~Gustafson,
T.~Xiao, S.~Whitehead, A.~C. Berg, W.~Lo, P.~Dollár, and R.~Girshick.
\newblock Segment anything.
\newblock In \emph{ICCV}, 2023.

\bibitem{li2019cosegnet}
Z.~Li, C.~Bi, J.~Zheng, J.~Jiang, and C.~Shi.
\newblock Co-segmentation inspired attention networks for video foreground
segmentation.
\newblock In \emph{ICCV}, 2019.

\bibitem{oquab2023dinov2}
M.~Oquab, T.~Darcet, T.~Moutakanni, H.~V. Vo, M.~Szafraniec, V.~Khalidov,
P.~Fernandez, D.~Haziza, F.~Massa, A.~El-Nouby, M.~Assran, N.~Ballas,
W.~Galuba, R.~Howes, P.-Y. Mairal, S.~Sherif, I.~Misra, M.~Bojanowski,
O.~Sivic, T.~Zisserman, J.~Moutakanni, I.~Misra, P.~Bojanowski, A.~Joulin,
A.~Vedaldi, and J.~Mairal.
\newblock DINOv2: Learning robust visual features without supervision.
\newblock \emph{TMLR}, 2023.

\bibitem{ren2024grounded}
T.~Ren, S.~Liu, A.~Zeng, J.~Lin, K.~Li, H.~Cao, J.~Chen, X.~Huang, Y.~Chen,
F.~Yan, Z.~Yang, H.~Zhang, C.~Li, J.~Liu, H.~Zhu, and L.~Zhang.
\newblock Grounded SAM: Assembling open-world models for diverse visual tasks.
\newblock \emph{arXiv:2401.14159}, 2024.

\bibitem{rother2004grabcut}
C.~Rother, V.~Kolmogorov, and A.~Blake.
\newblock GrabCut: Interactive foreground extraction using iterated graph cuts.
\newblock \emph{ACM ToG}, 23(3):309--314, 2004.

\bibitem{rother2006cosegmentation}
C.~Rother, T.~Minka, A.~Blake, and V.~Kolmogorov.
\newblock Cosegmentation of image pairs by histogram matching — incorporating
a global constraint into MRFs.
\newblock In \emph{CVPR}, 2006.

\bibitem{shotton2006textonboost}
J.~Shotton, J.~Winn, C.~Rother, and A.~Criminisi.
\newblock TextonBoost: Joint appearance, shape and context modeling for
multi-class object recognition and segmentation.
\newblock In \emph{ECCV}, 2006.

\bibitem{simonyan2015very}
K.~Simonyan and A.~Zisserman.
\newblock Very deep convolutional networks for large-scale image recognition.
\newblock In \emph{ICLR}, 2015.

\bibitem{vaswani2017attention}
A.~Vaswani, N.~Shazeer, N.~Parmar, J.~Uszkoreit, L.~Jones, A.~N. Gomez,
L.~Kaiser, and I.~Polosukhin.
\newblock Attention is all you need.
\newblock In \emph{NeurIPS}, 2017.

\bibitem{vicente2011object}
S.~Vicente, V.~Kolmogorov, and C.~Rother.
\newblock Object cut.
\newblock In \emph{CVPR}, 2011.

\bibitem{wu2023cycleseg}
Z.~Wu, J.~Su, M.~Tao, and T.~Wang.
\newblock CycleSegNet: Object co-segmentation with cycle refinement and
random walk.
\newblock \emph{TIP}, 2023.

\bibitem{zhang2015cosaliency}
K.~Zhang, T.~Chen, B.~Liu, and J.~Liu.
\newblock Co-saliency detection via a self-paced multiple-instance learning
framework.
\newblock \emph{TPAMI}, 2015.

\bibitem{zhang2020gradient}
L.~Zhang, J.~Dai, H.~Lu, Y.~He, and G.~Wang.
\newblock A bi-directional message passing model for salient object detection.
\newblock In \emph{CVPR}, 2018.

\end{thebibliography}
\endgroup

% ──────────────────────────────────────────────────────────────────────────────
\newpage
\onecolumn
\appendix
\section*{Appendix A: Full Source Code}
\label{app:code}

All code is available at \url{https://github.com/YOUR_USERNAME/SAMCo}.

\subsection*{A.1 cosegmentation.py — Full Pipeline}
\begin{verbatim}
class SAMCo:
    def __init__(self, sam_model_type, sam_checkpoint,
                 n_fg_points=5, n_bg_points=3,
                 top_k_ratio=0.30, n_refine_iter=2):
        self.feat_extractor = DINOv2FeatureExtractor()
        self.prompter       = SemanticConsensusPrompter(
            n_fg_points=n_fg_points, n_bg_points=n_bg_points,
            top_k_ratio=top_k_ratio)
        self.sam            = SAMWrapper(sam_model_type, sam_checkpoint)
        self.n_refine_iter  = n_refine_iter

    def segment(self, images):
        orig_sizes = [(img.height, img.width) for img in images]
        all_feats  = self.feat_extractor.extract_batch(images)
        saliency   = self.prompter.compute_consensus_saliency(all_feats)
        prompts    = self.prompter.generate_prompts(
                         saliency, self.feat_extractor, orig_sizes)
        results    = self.sam.predict_batch(images, prompts)
        masks      = [m for m, _ in results]
        for _ in range(self.n_refine_iter):
            prompts = self.prompter.refine_prompts_with_masks(
                          prompts, masks, saliency,
                          self.feat_extractor, orig_sizes)
            results = self.sam.predict_batch(images, prompts)
            masks   = [m for m, _ in results]
        return masks
\end{verbatim}

\subsection*{A.2 consensus\_prompting.py — SCP Core}
\begin{verbatim}
def compute_consensus_saliency(self, all_feats):
    N = len(all_feats)
    flat = [f.reshape(-1, f.shape[-1]) for f in all_feats]
    saliency_maps = []
    for i in range(N):
        scores = np.zeros(flat[i].shape[0])
        for j in range(N):
            if i == j: continue
            # cosine similarity: (P, P)
            Fi = flat[i] / (np.linalg.norm(flat[i], axis=1, keepdims=True) + 1e-8)
            Fj = flat[j] / (np.linalg.norm(flat[j], axis=1, keepdims=True) + 1e-8)
            sim = Fi @ Fj.T
            max_sim = sim.max(axis=1)
            k = max(1, int(self.top_k_ratio * len(max_sim)))
            threshold = np.partition(max_sim, -k)[-k]
            scores += np.where(max_sim >= threshold, max_sim, 0.0)
        scores /= (N - 1)
        smap = (scores - scores.min()) / (scores.max() - scores.min() + 1e-8)
        smap = smap.reshape(37, 37)
        saliency_maps.append(gaussian_filter(smap, sigma=1.0))
    return saliency_maps
\end{verbatim}

\subsection*{A.3 Dataset Links}
\begin{itemize}
  \item iCoseg: \url{https://chenlab.ece.cornell.edu/projects/touch-coseg/}
  \item MSRC v2: \url{https://www.microsoft.com/en-us/research/project/image-understanding/}
  \item Cosal2015: \url{http://www.zengwei.site/CoSal2015.html}
  \item SAM checkpoint: \url{https://dl.fbaipublicfiles.com/segment_anything/sam_vit_h_4b8939.pth}
  \item DINOv2 (via torch.hub): \url{https://github.com/facebookresearch/dinov2}
\end{itemize}

\section*{Appendix B: Individual Contributions}

This project is an individual submission.
All components — the SCP algorithm, ICR loop, Gradio demo, evaluation harness,
and this write-up — were implemented and written by Roll No.\ 230384.
The implementation builds upon the publicly released weights of SAM
(Meta AI~\cite{kirillov2023sam}) and DINOv2 (Meta AI~\cite{oquab2023dinov2});
no other external code was used beyond standard libraries (PyTorch, scikit-learn,
Pillow, Gradio).

\end{document}
